\section{Experiments}
\label{sec:experiments}

We evaluate \dgtdm on two motion-centric action recognition benchmarks,
Diving48 and Something-Something V2 (SSv2). Below we first describe the
datasets, evaluation protocol, and implementation details, then compare
\dgtdm with existing methods on the two benchmarks, and finally present
ablation studies on Diving48 that isolate the contribution of each branch.

\subsection{Setup}
\label{sec:setup}

\textbf{Datasets.}
We evaluate on two benchmarks. Diving48 \cite{diving48} contains approximately
$18$k trimmed video clips of $48$ dive classes; its videos share similar pool
backgrounds and differ in subtle body movements such as rotation angle and
timing. We use the official V2 train/test split. Something-Something V2
\cite{goyal2017ssv2} is a large-scale motion-centric benchmark of around $220$k
videos in $174$ classes of everyday object manipulations, such as pushing,
pulling, and dropping; the classes are defined by the direction and manner of
the motion rather than by the objects or the scene, so success requires
temporal reasoning over the whole sequence. For both datasets we report top-1
and top-5 accuracy.

\textbf{Implementation details.}
We use a ResNet-50 backbone as the temporal action encoder. The frozen spatial
branch is a CLIP ViT-B/16 encoder whose parameters are kept frozen. The DG-TDM
module is inserted after the $3\times3$ convolution of the last bottleneck in
each of the four stages, with a reduction ratio of $4$; the STG block is
inserted after each of the four stages, with a window size of
$7$. We train with stochastic gradient descent, momentum $0.9$, weight decay
$5\times10^{-4}$, and a step schedule. The classifier uses label-smoothed
cross-entropy with smoothing $0.1$. Training is performed on two NVIDIA
V100-SXM2-32GB GPUs. Our main setting samples $T=8$ frames per video. Full
hyperparameters are listed in the appendix.

\textbf{Baselines.}
We compare against representative methods that report results on Diving48 and
SSv2. Reported numbers of the compared methods are taken from their original
papers (with the source cited in the table).

\subsection{Comparison with state of the art}
\label{sec:main}

Tables~\ref{tab:diving48} and \ref{tab:ssv2} compare \dgtdm with existing
methods on the two benchmarks. We report values exactly as published in the
original papers, together with the input configuration (frames $\times$
resolution) used by each method. The compared methods span specialized
fine-grained models (DSTS, CAN), lightweight 2D temporal models (TANet,
TEINet, GST, AdaFuse, TDN, and AK-Net), and recent adaptation baselines
(MANet, LoKi, and AIM). For Diving48, the numbers under the DSTS protocol
are as reported by DSTS
\cite{li2022dsts} and those under the DVT protocol by DVT \cite{dvt2022}; the
remaining rows come from the original publications. For SSv2, all baseline values are taken from the cited original papers and
are restricted to configurations with an explicitly reported $8$-frame input.
We evaluate \dgtdm on both benchmarks with $8$ frames.

\begin{table}[t]
  \centering
  \caption{Comparison on Diving48. Metrics are top-1 and top-5 accuracy (\%).
  The Input / views column preserves the notation used by each source: an
  explicit $frames\times clips\times crops$ value is shown when reported,
  while sources that report only the number of input frames retain a single
  frame-count value. Resolution-based configurations are also kept in their
  original form. Top-5 values are marked ``-'' where the source does not
  report them. ``$\dagger$'' marks the DSTS protocol
  \cite{li2022dsts} and ``$\ddagger$'' the DVT protocol \cite{dvt2022}; the
  two papers use different evaluation protocols, so values across the two
  groups are not directly comparable. The TDN row reports the corrected
  official result compiled in \cite{chen2022aknet}, using the updated
  Diving48 annotations. The CAN row takes the $8$-frame ablation
  configuration from \cite{li2025can}, whose main result uses $16$ frames;
  the AIM row reports the $8$-frame result reproduced by \cite{cai2025loki}.}
  \label{tab:diving48}
  \small
  \begin{tabular}{llccc}
    \toprule
    Method & Backbone & Input / views & Top-1 (\%) & Top-5 (\%) \\
    \midrule
    \multicolumn{5}{l}{\emph{Protocol of DSTS \cite{li2022dsts}}} \\
    Video Swin-B \cite{liu2022videoswin} & Swin-B & - & 80.5$\dagger$ & - \\
    Swin-B + DSTS \cite{li2022dsts} & Swin-B & - & 83.0$\dagger$ & - \\
    \midrule
    \multicolumn{5}{l}{\emph{Protocol of DVT \cite{dvt2022}}} \\
    DVT \cite{dvt2022} & ViT-B & 16$\times$224$^2$ & 83.5$\ddagger$ & - \\
    \midrule
    \multicolumn{5}{l}{\emph{Methods with lightweight backbones}} \\
    TDN \cite{wang2021tdn} (corrected) & ResNet-50 & 8 & 80.5 & - \\
    AK-Net (TSM) \cite{chen2022aknet} & ResNet-50 & 8 & 79.3 & - \\
    AK-Net (TDN) \cite{chen2022aknet} & ResNet-50 & 8 & 81.7 & - \\
    CAN \cite{li2025can} (8-frame config) & ResNet-50 & 8$\times1\times1$ & 82.7 & 98.3 \\
    MANet \cite{li2025manet} & ResNet-50 & 8$\times1\times1$ & 81.8 & 98.0 \\
    LoKi \cite{cai2025loki} & ViT-B/16 & 8$\times1\times1$ & 78.4 & 98.2 \\
    AIM \cite{yang2023aim} & ViT-B/16 & 8$\times1\times1$ & 77.4 & 97.8 \\
    \midrule
    \multicolumn{5}{l}{\emph{Ours}} \\
    \dgtdm & ResNet-50 + frozen CLIP & 8$\times1\times1$ & \textbf{83.55} & \textbf{97.82} \\
    \bottomrule
  \end{tabular}
\end{table}

\begin{table}[t]
  \centering
  \caption{Comparison on Something-Something V2 with $8$ input frames and
  lightweight ResNet-50 backbones. Metrics are top-1 and top-5 accuracy (\%).
  TANet follows its original $8\times1\times1$ protocol \cite{liu2021tam};
  TDN and AK-Net follow the unified $8$-frame comparison in AK-Net Table II
  \cite{chen2022aknet}; ILA, STAN-conv, and M2-CLIP follow their reported
  multi-view 8-frame settings \cite{tu2023ila,liu2023stan,wang2024m2clip}.
  FEASE-Net, CAN, and MANet use their reported $8\times1\times1$
  configurations \cite{zhou2024fease,li2025can,li2025manet}. A dash denotes
  a top-5 value not reported by the source.}
  \label{tab:ssv2}
  \small
  \begin{tabular}{llccc}
    \toprule
    Method & Backbone & Frames & Top-1 (\%) & Top-5 (\%) \\
    \midrule
    \multicolumn{5}{l}{\emph{ResNet-50 and lightweight baselines}} \\
    TANet \cite{liu2021tam} (IN) & ResNet-50 & 8$\times1\times1$ & 60.5 & 86.2 \\
    TEINet \cite{liu2020teinet} (IN) & ResNet-50 & 8$\times1$ & 61.3 & - \\
    GST \cite{luo2019gst} (IN) & ResNet-50 & 8$\times1$ & 61.6 & - \\
    AdaFuse \cite{meng2021adafuse} (IN) & ResNet-50 & 8$\times1$ & 59.8 & - \\
    TDN \cite{wang2021tdn} (IN, corrected) & ResNet-50 & 8$\times1$ & 64.0 & - \\
    AK-Net (TSM) \cite{chen2022aknet} (IN) & ResNet-50 & 8$\times1$ & 60.2 & - \\
    AK-Net (TDN) \cite{chen2022aknet} (IN) & ResNet-50 & 8$\times1$ & 64.3 & - \\
    \midrule
    \multicolumn{5}{l}{\emph{Recent 8-frame adaptation baselines}} \\
    ILA \cite{tu2023ila} (ViT-B/16) & ViT-B/16 & 8$\times4\times3$ & 65.0 & 89.2 \\
    STAN-conv \cite{liu2023stan} (CLIP) & ViT-B/16 & 8$\times1\times3$ & 65.2 & 90.5 \\
    M2-CLIP \cite{wang2024m2clip} (CLIP) & ViT-B/16 & 8$\times1\times3$ & 66.9 & 90.1 \\
    FEASE-Net \cite{zhou2024fease} & ResNet-50 & 8$\times1\times1$ & 64.2 & 88.8 \\
    CAN \cite{li2025can} & ResNet-50 & 8$\times1\times1$ & 63.9 & 89.1 \\
    MANet \cite{li2025manet} & ResNet-50 & 8$\times1\times1$ & 63.6 & 88.8 \\
    \midrule
    \dgtdm (ours) & ResNet-50 + frozen CLIP & 8$\times1\times1$ & \textbf{65.3} & \textbf{89.1} \\
    \bottomrule
  \end{tabular}
\end{table}

Under the $8\times1\times1$ single-clip protocol, \dgtdm reaches
$83.55\%$/$97.82\%$ top-1/top-5 on Diving48 and $65.3\%$/$89.1\%$ on SSv2.
Its top-1 accuracy is the highest among the compared methods with a ResNet-50
backbone on both benchmarks.

\subsection{Ablation studies}
\label{sec:ablation}

We ablate \dgtdm on Diving48, the benchmark whose classes share the most
similar backgrounds and therefore stresses localization the most. All
ablations use $8$ frames. We organize the analysis into three parts: the
progressive build-up of the framework (Table~\ref{tab:ablation}), the
internal design choices of each branch (Table~\ref{tab:design}), and the
sensitivity to the main hyperparameters (Table~\ref{tab:hyper}).

\textbf{Progressive build-up.}
Table~\ref{tab:ablation} adds the components of \dgtdm to a TSN-style
baseline (a plain ResNet-50 that samples frames, extracts 2D features, and
classifies their temporal average), one component at a time, using the
$8\times1\times1$ protocol of the main comparison. Each step isolates one
mechanism: MoDeA contributes short-term cross-frame alignment, TSA adds a
persistent long-term state and thereby completes DG-TDM, and the STG branch
finally brings in the frozen-CLIP prior as guidance.

\begin{table}[t]
  \centering
  \caption{Progressive ablation on Diving48 with $8\times1\times1$ views.
  Each row adds one component of \dgtdm to the previous configuration; the
  last column reports the top-1 gain over the preceding row. Metrics are
  top-1 and top-5 accuracy (\%).}
  \label{tab:ablation}
  \small
  \begin{tabular}{lccc}
    \toprule
    Configuration & Top-1 (\%) & Top-5 (\%) & $\Delta$Top-1 \\
    \midrule
    TSN-style baseline (ResNet-50) & 71.36 & 95.82 & --- \\
    + MoDeA (short-term motion alignment) & 79.09 & 96.74 & +7.73 \\
    + TSA (long-term state memory; DG-TDM complete) & 81.57 & 97.77 & +2.48 \\
    + STG guidance (full \dgtdm) & \textbf{83.55} & \textbf{97.82} & +1.98 \\
    \bottomrule
  \end{tabular}
\end{table}

MoDeA raises top-1 by $7.73$ points over appearance averaging alone,
confirming that following the moving body part across frames is essential
when the discriminative signal is a small local displacement. TSA adds a
further $2.48$ points, showing that a state that aggregates the whole
sequence captures information that local alignment leaves out. The STG
branch contributes the final $1.98$ points, which measures the value of
bilateral interaction and adaptive gating under the same temporal encoder.
Top-5 accuracy improves consistently along the sequence, with the full model
reaching $97.82\%$.

\textbf{Design choices.}
Table~\ref{tab:design} switches off one internal design choice at a time
from the full model, so that each row keeps the affected component present
and isolates a single factor. The overall contribution of the MoDeA branch
is already measured by Table~\ref{tab:ablation} (removing it recovers the
$+7.73$ points of the baseline row), so the rows below address the choices
inside the remaining machinery. One factor is the fusion of the two DG-TDM
branches, which are combined multiplicatively so that TSA selects the
MoDeA-refined local features with a per-channel gate; we compare this with
adding the two branch outputs in parallel. A second factor is the direction
of the TSA scan, and we remove its backward pass. The bilateral nature of
the STG guidance is a third factor, and we keep only the prior-to-feature
direction. The last factor is where the frozen CLIP prior is applied: we
test injecting it at the last stage only, asking whether the position
knowledge must be used progressively at several spatial scales or once at
the top. The full model serves as the reference row.

\begin{table}[t]
  \centering
  \caption{Design-choice ablation on Diving48 with $8\times1\times1$ views.
  Each row switches off one internal design choice of the full \dgtdm while
  keeping the affected component and the remaining configuration fixed.
  Metrics are top-1 and top-5 accuracy (\%).}
  \label{tab:design}
  \small
  \begin{tabular}{lcc}
    \toprule
    Configuration & Top-1 (\%) & Top-5 (\%) \\
    \midrule
    Full \dgtdm (reference) & \textbf{83.55} & \textbf{97.82} \\
    DG-TDM: additive branch outputs (w/o TSA gate selection) & 83.41 & 98.02 \\
    TSA: forward-only scan (w/o backward context) & 82.58 & 97.76 \\
    STG: one-directional guidance (w/o $Y\!\to\!V$) & 82.70 & 97.88 \\
    CLIP prior: last-stage-only guidance (w/o per-stage) & 81.83 & 97.65 \\
    \bottomrule
  \end{tabular}
\end{table}

\textbf{Hyperparameter sensitivity.}
Table~\ref{tab:hyper} verifies that the reported gain of the guidance is
robust to the two settings that define how the STG block injects the frozen
CLIP prior into the recognition branch: the window size $w$ and the residual
bound $\lambda_{\max}$, swept around the defaults of the full model
($w{=}7$, $\lambda_{\max}{=}0.2$; last row of Table~\ref{tab:ablation}).
The window size fixes the spatial extent of the bilateral attention: too
small a window cannot cover the moving action region, while too large a
window lets background context dilute the prior. The residual bound caps how
strongly a stage may be modified: a low bound starves the injection, and a
high bound lets the prior overwrite the appearance features of the lower
branch. Because the residual scale is itself learned within
$[0,\lambda_{\max}]$, the bound only sets the allowed headroom. Across both
sweeps the default configuration stays on a plateau in top-1: the window
$w{=}7$ outperforms both of its neighbors, and tightening the residual bound
to $0.1$ is the only change that costs a measurable amount, while loosening
it to $0.4$ brings no further gain, confirming that the guidance does not
depend on a fine-tuned setting.

\begin{table}[t]
  \centering
  \caption{Hyperparameter sensitivity of the STG guidance on Diving48 with
  $8$ frames. The starred value marks the default of the full model
  (Table~\ref{tab:ablation}); each hyperparameter is swept to one smaller
  and one larger value. Metrics are top-1 and top-5 accuracy (\%).}
  \label{tab:hyper}
  \small
  \begin{tabular}{llcc}
    \toprule
    Hyperparameter & Value & Top-1 (\%) & Top-5 (\%) \\
    \midrule
    STG window size $w$ & 5 & 83.29 & 97.74 \\
     & $\mathbf{7}^{\ast}$ & \textbf{83.55} & \textbf{97.82} \\
     & 9 & 83.14 & 97.82 \\
    \midrule
    Residual bound $\lambda_{\max}$ & 0.1 & 82.95 & 97.38 \\
     & $\mathbf{0.2}^{\ast}$ & \textbf{83.55} & \textbf{97.82} \\
     & 0.4 & 83.37 & 98.05 \\
    \bottomrule
  \end{tabular}
\end{table}

\subsection{Efficiency analysis}
\label{sec:efficiency}

Because the frozen CLIP branch does not back-propagate and the DG-TDM operates
in a reduced bottleneck dimension, the framework adds modest trainable
parameters and FLOPs over the temporal action encoder baseline. The appendix
reports the parameter count, FLOPs, and per-clip throughput of the main
variants.

\subsection{Qualitative analysis}
\label{sec:qual}

The learned mechanisms are inspected on sampled Diving48 frames in two ways.
The deformation offsets predicted by MoDeA show whether different channel
groups track different body-part motions, and the gates produced by the STG
branch show whether the lower branch focuses on the action region instead of
the shared background.
